\subsection{Класифікація візуальних артефактів масштабування}

Будь-який процес геометричного перетворення та просторового масштабування цифрових зображень супроводжується виникненням специфічних візуальних спотворень. Їхня наявність зумовлена математичними обмеженнями алгоритмів інтерполяції. Згідно з проведеними дослідженнями, основними типами таких інтерполяційних артефактів є:

\begin{itemize}
	\item \textbf{Ефект дзвону (Ringing):} Поява одиночної паразитної хвилі або контуру безпосередньо поблизу різких меж на зображенні.
	
	\item \textbf{Перерегулювання (Overshooting):} Виникнення цілої серії подібних хвиль навколо контрастних переходів.
	
	\item \textbf{Субпіксельний зсув (Sub-pixel shift):} Небажана загальна трансляція (зміщення) всього зображення у певному просторовому напрямку. Це явище зазвичай є наслідком специфічних особливостей програмної реалізації конкретного алгоритму. Хоча візуально такий зсув може бути абсолютно непомітним для людського ока і не впливати на суб'єктивну якість, він має критично негативний вплив на формальні об'єктивні метрики (наприклад, суттєво погіршує значення PSNR та MSE при попіксельному порівнянні з еталоном).
	
	\item \textbf{Аліасинг (Aliasing):} Так званий «східчастий ефект» (step effect), який проявляється як нерівномірність та «зубчастість» різких діагональних контурів. Цей дефект виникає як при обробці синтетичних (векторних), так і реальних фотографічних зображень. Найбільш характерним аліасинг є для алгоритмів типу «найближчий сусід».
	
	\item \textbf{Розмиття (Blurring):} Втрата чіткості та загальне зниження різкості зображення після масштабування. Більшою чи меншою мірою цей ефект притаманний усім інтерполяційним алгоритмам, але найсильніше він виражений у простих методах (таких як білінійна інтерполяція).
\end{itemize}

В теорії цифрової обробки зображень існує фундаментальний компроміс: штучне підвищення різкості алгоритмічним шляхом призводить до посилення інших артефактів (аліасингу, перерегулювання), і навпаки — спроби придушити артефакти неминуче викликають ще більше розмиття зображення. Для часткової компенсації цих негативних ефектів на практиці застосовуються додаткові процедури просторової фільтрації. Використання різних комбінацій низькочастотних та високочастотних фільтрів дозволяє згладити або, навпаки, підкреслити певні деталі, покращуючи загальне візуальне сприйняття результату.